\section{Class presentations and reductions}
\label{sec:classes}

\subsection{First-order definitions of P, NP, and their relatives}

In every language-equivalence and reduction statement, input variables range
over valid word codes \(x\ge1\). We write this guard explicitly in the
defining sentences below and leave it implicit in subsequent displays.
Transducers output valid word codes, including code 1 for the empty word.
The reserved code 0 is outside the language universe; in particular,
complementation is relative to valid words.

A deterministic polynomial descriptor is \(d=(e,c,k)\), with \(c\ge1\).
Malformed descriptors are normalized to a fixed rejecting descriptor;
the raw description component \(e=0\) is normalized as specified above.
Define the Boolean function
\[
 P_d(x)=\Acc(e,x,p_{c,k}(\len x)).
\]
A verifier descriptor \(v=(e,c,k,a,l)\) has witness length
\(p(n)=c(n+1)^k\) and time clock \(q(n)=a(n+1)^l\). Its input record is
\(x\#y\), where \(y\) has exactly \(p(n)\) bits and the separator is part
of the fixed finite machine alphabet. To keep binary word codes and this
initialized tape distinct, define \(\Acc_{\#}(e,x,y,t)\) by the same step
iteration as \(\Acc\), but starting with the tape record \(x\#y\).
This is a fixed primitive-recursive function of two binary word codes and
the clock. Its simulation cost uses tape-input length \(\len x+\len y+1\).
Equivalently, an ordinary binary machine can read a self-delimiting binary
pair and construct this record with polynomial overhead.
Let \(\operatorname{bits}_p(j)\)
denote the \(p\)-bit representation of \(j<2^p\), including leading zeros.
Set
\[
 N_v(x)=\bigvee_{j<2^{p(n)}}
 \Acc_{\#}(e,x,\operatorname{bits}_{p(n)}(j),q(n)),\qquad n=\len x.
\]
All operations in this definition are fixed primitive-recursive functions.
An exact-length certificate entails no loss: a certificate of length at most
\(p\) can be represented by a unary length delimiter, its contents, and
padding in a fixed length \(2p+1\) record.

\begin{definition}[Representing a language]
For a fixed arithmetic formula \(A(x)\), write
\[
 \operatorname{PRep}(d,A):=\forall x\ge1\bigl(A(x)\leftrightarrow P_d(x)=1\bigr),
\]
and define \(\operatorname{NPRep}(v,A)\) by replacing \(P_d\) with \(N_v\).
The abbreviations \(A\in\Pclass\), \(A\in\NP\) mean existence of the
corresponding numerical descriptor. This use of \(A\) is schematic; there is
no internal quantification over arbitrary predicates. Alternatively, work
entirely with the Boolean functions \(P_d,N_v\), and compare them by
\(\forall x\ P_d(x)=N_v(x)\).
\end{definition}

Exponential time uses \(\Acc(e,x,e_{c,k}(\len x))\). Space presentations
use acceptance in the capped configuration graph for an explicit
\(s(n)=c\lceil\log_2(n+2)\rceil\) or \(s(n)=p_{c,k}(n)\), as developed
below. Integer logarithms mean bounded search for a power of two, not real
analysis. Deterministic versus nondeterministic edges distinguish
\(\Lclass,\NL\) and \(\PSPACE,\NPSPACE\).

For coNP use \(1-N_v(x)\). Its direct logical form is a bounded universal
statement that every certificate is rejected. Each such input statement is
decidable in \ipra. This does not place its language in NP, because the
exhaustive decision algorithm can take exponential time.

\begin{theorem}[Verifier and branch presentations]
Polynomial-time nondeterministic acceptance and the verifier presentation
\(N_v\) agree, with primitive-recursive translations of descriptors and
bounded witness/branch correctness proofs in \ipra.
\end{theorem}

\begin{proof}
Normalize a fixed finitely branching transition table so that each choice
is coded by a fixed number \(b\) of bits. A branch of length at most \(T\)
is coded by \(bT\) choice bits, with padding after halting. A deterministic
verifier replays those choices, rejecting an invalid transition code. Its
simulation has polynomial overhead by the preceding sweep implementation.
The witness constructor records the choices of a supplied branch, and the
inverse replays them. The prefix invariant asserts equality of the recorded
and replayed configurations at every earlier time, a bounded formula.

Conversely, a nondeterministic machine guesses the prescribed polynomial
number of certificate bits, then runs the fixed clocked verifier. Its clock
includes guessing, input assembly and verification. The constructed branch
exists exactly when the bounded witness disjunction is 1. All loops and
record lengths are explicitly bounded. No infinite branch principle is used.
\end{proof}

\subsection{Elementary closure and exhaustive simulation}

\begin{theorem}[Closure and basic containments]
\label{thm:closure}
\ipra{} proves the numerical compiler theorems corresponding to:
\(\Pclass\subseteq\NP\); closure of P under complement, finite union and
intersection; closure of NP under finite union and intersection; and
\(\NP\subseteq\EXP\).
\end{theorem}

\begin{proof}
For \(\Pclass\subseteq\NP\), ignore the witness and run the supplied
deterministic descriptor. To complement P, evaluate the entire clocked
Boolean computation and negate its answer, including the timeout answer.
Merely interchanging the original machine's halting labels is insufficient
when the original machine can time out. Sequentially evaluate two clocked
descriptors and combine their bits by AND or OR to obtain intersection or
union. The resulting sums and simulation overhead have polynomial majorants.

For intersection of two NP presentations, concatenate their two
certificates and verify both. The forward witness map concatenates; the
reverse projects. For union use a selector bit and a certificate field
padded to the larger of the two prescribed lengths; the selector determines
which verifier is run. These are explicit witness maps, so the existential
equivalences are intuitionistic.
If a concatenated, padded, or selector-augmented certificate length is not
already a monomial, choose the explicit monomial majorant from
Section~\ref{sec:machines} and ignore the additional padding. This keeps the
constructor inside the displayed descriptor family.

Finally enumerate \(2^{p(n)}\) certificates, replaying the verifier for
\(q(n)\) steps on each. A fixed implementation takes at most
\[
 A\,2^{p(n)}(n+p(n)+q(n)+1)^b
\]
steps for fixed integers \(A,b\). Replace the polynomial factor and \(A\)
by powers of two using \(r\le2^r\), and absorb the resulting polynomial
exponent into \(2^{r_0(n)}\). The scan's OR accumulator is exactly the
defining bounded disjunction. Its induction invariant compares the
accumulator with the finite prefix OR. Exhaustive search here proves an
exponential upper bound, not a polynomial one.
\end{proof}

\subsection{Reductions carry algorithms and bounds}

\begin{definition}[Polynomial many-one reduction]
A reduction from Boolean presentations \(A\) to \(B\) is a clocked
polynomial transducer \(f\), together with
\[
 \forall x\quad A(x)=B(f(x)).
\]
For formula presentations replace equality by equivalence. Its certificate
records a polynomial time majorant and an output-length majorant. A
\emph{witness-preserving} reduction additionally specifies primitive-recursive
maps transporting source witnesses to target witnesses and target witnesses
back to source witnesses, with their correctness and size bounds.
\end{definition}

\begin{theorem}[Composition and downward closure]
\label{thm:reduction}
Polynomial many-one reduction is reflexive and transitive. If \(A\le_m^pB\)
and \(B\in\Pclass\), then \(A\in\Pclass\). If \(B\in\NP\), then
\(A\in\NP\). These statements have explicit descriptor compilers and
\ipra{} proofs.
\end{theorem}

\begin{proof}
The identity transducer gives reflexivity. For composition, run \(f\), then
\(g\) on its output. If the respective bounds are \(p(n)\) and \(q(m)\),
then \(\len{f(x)}\le p(n)\), so the useful work is at most
\(p(n)+q(p(n))\), with an explicit copying and simulation overhead.
The polynomial-majorant calculation in Section~\ref{sec:machines} supplies
a resulting clock. Correctness is ordinary equality transitivity:
\(A(x)=B(f(x))=C(g(f(x)))\).

For downward closure of P, compose with the target decider. For NP, compute
\(f(x)\), take a certificate for the target verifier at that input, and
verify it. If target witness length is \(r(m)\), use a fixed padded field
of length \(r(p(n))\), discarding padding after the required length
\(r(\len{f(x)})\). Its existence is equivalent by extension and projection.
The combined verification time is polynomial. No unbounded choice of a
machine or a witness occurs in these compiler proofs.
\end{proof}

A language is hard for a presentation class if every descriptor in the class
has such a reduction to it; completeness adds membership. In arithmetic this
can be written as a numerical quantification over descriptors, reduction
codes, and their pointwise correctness. The standard completeness proofs
below do more: they give a primitive-recursive constructor for the reduction
code rather than only asserting its existence.

\subsection{Search, functions, promises, and oracle queries}

The class FP is presented by polynomially clocked transducers. A search
problem has a Boolean polynomial-time relation \(R(x,y)\) and a polynomial
balance bound \(\len y\le p(\len x)\). A solution procedure returns a witness
when one exists, and a distinguished failure record otherwise. Total search
adds a proved assertion that a witness always exists; this assertion must
not be inferred just from the verifier's totality.

A promise problem is specified by a promise formula \(H(x)\), a yes-condition
and a no-condition, with the claimed algorithmic behavior required only under
\(H(x)\). Completeness arguments must preserve the promise, not just map
yes-instances correctly. A numerical certificate for a promise can be handled
as ordinary input data; an arbitrary truth oracle is not silently inserted.

A polynomial oracle reduction consists of a clocked query machine with
polynomially many polynomial-length queries. For a fixed Boolean oracle
presentation, its query answers can be substituted to obtain a primitive-
recursive bounded computation. If no primitive-recursive oracle definition
is supplied, use a finite transcript of queries and supplied answers and a
relational correctness statement. This gives an appropriate base for later
oracle complexity while respecting first-order arithmetic. No unrestricted
oracle truth predicate is part of the present theory.
